\section{Evaluation Infrastructure}
\sectionkicker{EVALUATION}
\begin{wideflow}
{\Large\bfseries 量る台――AgentPerfの再走、OpenTTSの物差し、RLの置き場}\par
\smallskip
{\color{SurveyMuted}台と模型を取り違えない。何が量れるかの紙であり、何ができたかの証しではない。}\par
\end{wideflow}

AA-AgentPerf-Localは、手元で走らせる再走の台である\cite{agentperf}。本物の軌跡をSparkや5090やHaloやM5Proでなぞり直す。構えは14通りすべてが仮のものと明記され、売り手の計測として置く。運びの話に限る品であり、使い手の賢さの話ではない。帯域の屋根線との比べは、上限の線に対する成功群の嵩上げ（Qwen3.8-27Bで3割から120\%ほど、構えによる幅）であって、素の対照ではない。一人で走らせる手元と、混み合う本番（メガワットあたりの使い手という別の物差し、6月12日の別の紙）は方法も指標も別物であり、数の貸し借りはしない。

Open TTS Leaderboardは、声の品を比べる台であり、新しい声の模型ではない\cite{opentts}。英語の巨視の平均の誤り、多言語の字の誤り、RTFxとTTFAとSIMを、Seed-TTSとCV3の上で量る。TTFAの手順は一句ずつ確かめた。同じ英語の五十の促し、一つずつの束ね、最初の三走を慣らしに除けて残りの中央値、H200の眺めが既定で、CPUの小さな伸びる集合も添う。人の好みは量っていないと紙が明記する。比べは写しの時点の話であり、普遍の番付にはしない。

時の置き場は欠かさない。場と物差しの本体は今週にあるが、評価の写しは9月30日の記事では「やがて開ける」という未来形である。置き場の初の写しの固まりは10月7日14時53分41秒（協定時）と確かめた。窓の締め切り（10月2日22時協定時）の時点では置き場に写しはなく、9月15日の置き場の生まれをもって公開と読んではならない。調べた置き場の範囲の話であり、どこにも出ていないことの証しではない。

RL Environments Hubは、強化学習の置き場の目録である\cite{rlenv}。rl-environmentの絞りと四つの札（harbor・verifiers・openenv・nemo-gym）、枠ごとの写しが並ぶ。束ねの資料と走りの報いは分け、新しい置き場の型はなく、札は相性の宣言であって変換や起動の話ではない。枠の版留め（harbor==0.21.0など）が束ね全体を留めるかのように読んではならない。束ねの版や選びの基準、走りの土台の固定は開いたままとして明記する。種の留めも、留まらない所は留まらないと書く。

\begin{claimboundary}[資料と評価の境界]
台の話と模型の話は別であり、運びの話を使い手の賢さに読み替えない。屋根線の比べは上限の線への嵩上げであり、素の対照ではない。写しの到着と置き場の生まれは別の事柄であり、未来形を公開済みと読まない。枠の版留めを束ね全体の留めと読まない。
\end{claimboundary}
